\documentclass[letterpaper,11pt]{moderncv}
\moderncvstyle{banking}
\moderncvcolor{black}

\usepackage[letterpaper, margin=0.8in]{geometry}
\usepackage{xcolor}
\usepackage{fontspec}
\setmainfont{Times New Roman}
\usepackage[unicode]{hyperref}

\name{Saeyoung}{Kim}
\address{Cambridge, MA, USA}{}{}
\phone[mobile]{+1 (351) 322 5510}
\email{saeyoung\_kim@uml.edu}
\social[linkedin]{saeyoung-macx-kim}

\recipient{Hiring Team}{Neros Technologies\\El Segundo, CA}
\date{\today}
\opening{Dear Hiring Team,}
\closing{Sincerely,}

\begin{document}

\makelettertitle

I'm applying for the Manufacturing Hardware Test Engineer role. I've spent the last several years designing test hardware, bringing up electromechanical systems, and building the infrastructure that makes repeatable testing possible. This role is basically the same problem at drone scale, and I'd like to be part of solving it.

At Hanwha Systems, I designed EGSE architecture for military satellite programs and owned the test side of integration from subsystem validation through flight model acceptance. That meant configuring modular test setups, commissioning data acquisition hardware, and writing the procedures that technicians ran on the floor. When things failed during integration, I did the root cause analysis, traced it back to whether it was a board issue, a wiring problem, or a design gap, and worked with the hardware teams to fix it. I also managed purchasing and integration of functional test systems, and pushed design-for-test feedback back into the next build.

My PhD was where I learned to build test infrastructure from nothing. I designed a multi-modal sensor suite for a CubeSat payload, which meant circuit design, PCB layout, cleanroom fabrication, firmware development in C/C++ talking over UART, SPI, and I2C, and then building custom test fixtures with oscilloscopes and DMMs to characterize everything. I went through multiple prototype-to-qualification cycles, running environmental and reliability tests each round to find failure modes and close them out.

More recently at UMass Lowell, I've been writing Python tools for automated process monitoring and anomaly detection in manufacturing. The work is focused on replacing commercial platforms with custom diagnostic software that's faster to iterate on and easier to adapt to new processes.

On the export compliance side: my green card is currently pending, and I expect to have it in hand within the next few months. I'm happy to discuss timing and work authorization details directly. I'm also willing to relocate to El Segundo for this role. Looking forward to talking more about how I can help build out your test capabilities.

\makeletterclosing

\end{document}
